\section{Butun olam tortishish qonuni}



\begin{enumerate}

\item Qaysi olim gravitatsiya doimiysini son qiymatini hisoblab chiqargan?
\begin{enumerate}[label=\Alph*)]
\item A. Eynshteyn
\item A. Avogadro
\item D. Mendeleyev
\item Kavendish
\end{enumerate}

\item Yer massasi Quyosh massasidan 330000 marta kichik. Quyosh tomonidan yerga ta'sir etuvchi tortish kuchining Yer tomonidan Quyoshga ta'sir etuvchi kuchga nisbatini toping?
\begin{enumerate}[label=\Alph*)]
\item 1
\item 1/575
\item 330000
\item 1/330000
\end{enumerate}

\item Ikki jismdan har birining massasi 2 marta oshirilsa va ular orasidagi masofa 2 marta kamaytirilsa, ularning o‘zaro tortishish kuchi qanday o'zgaradi?
\begin{enumerate}[label=\Alph*)]
\item 4 marta ortadi
\item 4 marta kamayadi
\item 8 marta kamayadi
\item 16 marta ortadi
\end{enumerate}

\item Ikki jism orasidagi gravitatsiya kuchi 64 marta kamaygan bo‘lsa, ular orasidagi masofa qanday o‘zgaradi?
\begin{enumerate}[label=\Alph*)]
\item o‘zgarmaydi
\item 8 marta ortadi
\item 8 marta kamayadi
\item 64 marta ortadi
\end{enumerate}

\item Jismning inert massasi va gravitatsion massasi qanday farq qiladi?
\begin{enumerate}[label=\Alph*)]
\item inert massa nolga teng
\item inert massa katta
\item gravitatsion massa nolga teng
\item farq qilmaydi
\end{enumerate}

\item Oyning massasi yerning massasidan 81 marta kichik. Yer va oy markazlarini birlashtiruvchi to‘g‘ri chiziqdagi Yerdan qanday masofadagi nuqtada jismning yerga tortilish kuchi va Oyga tortilish kuchi o‘zaro muvozanatlashadi?
\begin{enumerate}[label=\Alph*)]
\item 0,2 qismiga teng
\item 0,3 qismiga teng
\item 0,9 qismiga teng
\item 0,125 qismiga teng
\end{enumerate}

\item Massalari 1200 kg dan bo‘lgan ikki avtomobil massa markazlari orasidagi masofa 5 m ga teng. Avtomobillar bir-birini qanday kuch (N) bilan tortishadi?
\begin{enumerate}[label=\Alph*)]
\item $384 \cdot 10^{-7}$
\item $3,84 \cdot 10^{-7}$
\item $38,4 \cdot 10^{-7}$
\item $38,4 \cdot 10^{-5}$
\end{enumerate}

\item Yerning massasi $6,0 \cdot 10^{24}$ kg, Oyning massasi $7,35 \cdot 10^{22}$ kg, ularning markazlari orasidagi masofa 384400 km. Yer bilan Oy orasidagi tortishish kuchi topilsin (N).
\begin{enumerate}[label=\Alph*)]
\item $8,0 \cdot 10^{20}$
\item $2,0 \cdot 10^{18}$
\item $3,0 \cdot 10^{20}$
\item $2,0 \cdot 10^{20}$
\end{enumerate}

\item Kosmik raketa yer sathidan besh yer radiusiga teng bo‘lgan masofaga uzoqlashganda kosmik raketaning yerga tortishish kuchi necha marta kamayadi?
\begin{enumerate}[label=\Alph*)]
\item 6 marta
\item 36 marta
\item 9 marta
\item 81 marta
\end{enumerate}

\item Massalari $m$, $2m$ va $3m$ bo‘lgan uchta jism o‘rtasidagi tortishish kuchi qaysi juftlikda eng kichik? (1-2 masofa 2R, 2-3 masofa 3R)
\begin{enumerate}[label=\Alph*)]
\item 1 va 2
\item 2 va 3
\item 3 va 1
\item 1 va 3; 2 va 3
\end{enumerate}

\item Ikkita bir xil shar orasidagi tortishish kuchi 0,01 N. Agar ularning markazlari orasidagi masofa 10 m bo‘lsa, sharlarning massasi (kg) qanchaga teng?
\begin{enumerate}[label=\Alph*)]
\item $8 \cdot 10^{5}$
\item $1,2 \cdot 10^{5}$
\item $1,2 \cdot 10^{4}$
\item 60
\end{enumerate}

\item Yer sirtidan Yer radiusiga teng balandlikda erkin tushish tezlanishi nimaga teng ($m/s^2$). Yer sirtida $g = 10 m/s^2$.
\begin{enumerate}[label=\Alph*)]
\item 1,25
\item 10
\item 2,5
\item 5
\end{enumerate}

\item Jism osilgan dinamometr Yer sirtidan 1 N ni ko‘rsatadi. Agar jism yer sirtidan ikki Yer radiusiga teng masofaga chiqarilsa, Yerga necha nyuton kuch bilan tortiladi?
\begin{enumerate}[label=\Alph*)]
\item 3
\item 1/3
\item 1/9
\item 1/4
\end{enumerate}

\item Yer sirtidan qanday (km) balandlikda og‘irlik kuchi Yer sirtidagi og‘irlik kuchining 64 foizini tashkil qiladi. Yerning radiusi 6400 km.
\begin{enumerate}[label=\Alph*)]
\item 1400
\item 1600
\item 1500
\item 1700
\end{enumerate}

\item Agar Quyoshning radiusi yer radiusidan 108 marta katta, zichligi esa yer zichligidan 4 marta kichik bo‘lsa, Quyosh sirtida og‘irlik kuchining tezlanishi qanday bo‘ladi ($m/s^2$)?
\begin{enumerate}[label=\Alph*)]
\item 10,8
\item 27
\item 67,5
\item 270
\end{enumerate}

\item Yer sirtida kosmonavtga 720 N gravitatsion kuch ta’sir qilmoqda. Kosmonavt yer sirtidan ikki radiusga teng masofada bo'lganda unga qanday gravitatsion kuch ta’sir qiladi?
\begin{enumerate}[label=\Alph*)]
\item 180
\item 240
\item 80
\item 360
\end{enumerate}

\item Oyning Yer atrofida orbita bo‘ylab aylanishidagi markazdan qochma tezlanishi nimaga teng ($m/s^2$). Masofa 60 Yer radiusiga teng.
\begin{enumerate}[label=\Alph*)]
\item 1/360
\item 1/3600
\item 1/720
\item 1/60
\end{enumerate}

\item Yer sirtidan qanday masofada kosmik kemaning Yerga tortilish kuchi Yer sirtidagiga nisbatan 25 marta kichik bo‘ladi? (Yer radiuslarida)
\begin{enumerate}[label=\Alph*)]
\item $R$
\item $4 R$
\item $9 R$
\item $3 R$
\end{enumerate}

\item Jism osilgan dinamometr yer sirtida 1 N ni ko‘rsatadi. Agar jism yer sirtidan to‘rt yer radiusiga teng masofaga chiqarilsa, yerga necha nyuton kuch bilan tortiladi?
\begin{enumerate}[label=\Alph*)]
\item 3
\item 1/9
\item 1/4
\item 1/25
\end{enumerate}

\item Yer sirtidan qanday balandlikda (km) jismning og'irlik kuchi 4 marta kamayadi? $R = 6400$ km.
\begin{enumerate}[label=\Alph*)]
\item 19200
\item 25652
\item 12800
\item 6400
\end{enumerate}

\end{enumerate}


\begin{table}[h!]
\centering
\caption{17-\S. Butun olam tortish qonuni bo'yicha kalitlar}
\begin{tabular}{|l|c|c|c|c|c|c|c|c|c|c|}
\hline
\textbf{Test №} & 1 & 2 & 3 & 4 & 5 & 6 (7) & 7 (8) & 8 (9) & 9 (12) & 10 (13) \\ \hline
\textbf{Javob}  & D & A & D & B & D & C & C & D & B & C \\ \hline
\hline
\textbf{Test №} & 11 (15) & 12 (18) & 13 (19) & 14 (20) & 15 (22) & 16 (32) & 17 (33) & 18 (34) & 19 (35) & 20 (36) \\ \hline
\textbf{Javob}  & B & C & C & B & D & C & A & B & D & D \\ \hline
\end{tabular}
\end{table}










\section{Jismning og`irligi}



\begin{enumerate}

\item Agar jism ... ta’sirida harakatlansa, u vaznsizlik holatida bo‘ladi.
\begin{enumerate}[label=\Alph*)]
\item faqat og‘irlik kuchi
\item faqat elastiklik kuchi
\item markazga intilma kuch
\item harakatlanmasa
\end{enumerate}

\item Quyidagi jumlaning mazmuniga mos keladigan so‘z yoki birikmani tanlab, to‘liq gap hosil qiling: ... vaznsiz holati hosil bo‘ladi.
\begin{enumerate}[label=\Alph*)]
\item erkin tushishda
\item tinchlik holatida
\item to‘g‘ri chiziqli tekis harakatda
\item muvozanatda
\end{enumerate}

\item Arqon ko‘pi bilan $m$ massali yukni ko‘tara oladi. Shu arqon yordamida qanday massali yukni $g/2$ tezlanish bilan ko‘tarish mumkin?
\begin{enumerate}[label=\Alph*)]
\item $m/2$
\item $3m/4$
\item $m/3$
\item $2m/3$
\end{enumerate}

\item 50 kg massali havo shari 5 m/s tezlik bilan ko‘tarilmoqda va $5 m/s^2$ tezlanish bilan tezligi kamaymoqda. Uning vazni (N) qanday? $g=10 m/s^2$.
\begin{enumerate}[label=\Alph*)]
\item 250
\item 650
\item 500
\item 750
\end{enumerate}

\item Ko‘tarish mashinasining kanatga osilgan kleti yuqoriga yo‘nalgan $1,6 m/s^2$ tezlanish bilan harakat qilayotgan bo‘lsa, kanatning tarangligini aniqlang (N). Kletning massasi 300kg.
\begin{enumerate}[label=\Alph*)]
\item 4250
\item 7520
\item 3420
\item 1280
\end{enumerate}

\item Lift pastga tushayapti. Uning $3 m/s^2$ ga teng tezlanishi yuqoriga yo‘nalgan. Liftda 70 kg massali odam bor. Uning vazni (og‘irligi) necha nyuton?
\begin{enumerate}[label=\Alph*)]
\item 210
\item 490
\item 700
\item 910
\end{enumerate}

\item Ipga bog‘langan, massasi 2 kg bo‘lgan jism $5 m/s^2$ tezlanish bilan ko‘tarilmoqda. Ipning taranglik kuchini toping (nyutonlarda). $g=10 m/s^2$.
\begin{enumerate}[label=\Alph*)]
\item 5
\item 10
\item 20
\item 30
\end{enumerate}

\item Jismning og‘irlik kuchi havoning o‘rtacha qarshilik kuchiga teng bo‘lganda, jism 30 m balandlikdan necha sekundda tushadi? Jismning yerga urilish vaqtidagi tezligi 10 m/s.
\begin{enumerate}[label=\Alph*)]
\item 1,5
\item $\sqrt{6}$
\item 3
\item 1
\end{enumerate}

\item Ko‘tarish krani massasi 1 t bo‘lgan yukni ko‘tarmoqda. Agar yuk $25 m/s^2$ tezlanish bilan (juda qisqa vaqt) harakatlansa, ko‘tarilishning boshida trosning taranglik kuchi qancha bo‘lgan (kN)?
\begin{enumerate}[label=\Alph*)]
\item 30
\item 35
\item 70
\item 45
\end{enumerate}

\item Parashyut to‘la ochilgan paytda havoning parashyutchi va parashyutga ta’sir qiluvchi minimal qarshilik kuchi nimaga teng (N)? Parashyutchining parashyut bilan birgalikdagi massasi 75 kg.
\begin{enumerate}[label=\Alph*)]
\item 75
\item 735
\item 150
\item 300
\end{enumerate}

\item Tinch holatdagi og‘irligi 10 N bo‘lgan jismni vertikal yuqoriga ko‘tarish lozim. Yuk bog‘langan ip 20 N kuchga bardosh bera oladi. Shu ipdan foydalanilganda yukni 2 s davomida qanday maksimal balandlikka ko‘tarish mumkin (m)?
\begin{enumerate}[label=\Alph*)]
\item 19,6
\item 36,2
\item 28,4
\item 6,4
\end{enumerate}

\item Shaxta ko‘targichi platformasining og‘irligi 2,5 kN. Arqonning taranglik kuchi 2,0 kN bo‘lsa, ko‘targich qanday tezlanish bilan tushmoqda ($m/s^2$)?
\begin{enumerate}[label=\Alph*)]
\item 4
\item 2
\item 0,5
\item 8
\end{enumerate}

\item Lift kabinasining shiftiga mahkamlangan prujinali taroziga $m$ massali jism osilgan. Lift kabinasi yuqoriga yoki pastga tekis harakatlanganda, prujinali tarozining ko‘rsatishi $F$ qanday bo‘ladi?
\begin{enumerate}[label=\Alph*)]
\item $F = mg$
\item $F > mg$
\item $F < mg$
\item $F = 0$
\end{enumerate}

\item Pastga tushayotgan lift kabinasi polida yotgan 140 kg massali yuk polga 1440 N kuch bilan bosadi. Lift tezlanishining kattaligini ($m/s^2$) va yo‘nalishini aniqlang.
\begin{enumerate}[label=\Alph*)]
\item 0,48, liftning tezlanishi pastga yo‘nalgan
\item 0,48, liftning tezlanishi yuqoriga yo‘nalgan
\item 10,5, liftning tezlanishi yuqoriga yo‘nalgan
\item 10,5, liftning tezlanishi pastga yo‘nalgan
\end{enumerate}

\item Raketa gorizontal yo‘nalishda $2g$ tezlanish bilan harakatlanmoqda. Undagi 2 kg massali jism yuklanishi qanday? $g$ — erkin tushish tezlanishi.
\begin{enumerate}[label=\Alph*)]
\item $\sqrt{5}$
\item 2
\item $\sqrt{3}$
\item 3
\end{enumerate}

\item Raketa yer sirtiga nisbatan $30^\circ$ burchak ostida $2g$ tezlanish bilan ko‘tarilmoqda. Raketadagi $m$ massali kosmonavtning vazni qanday?
\begin{enumerate}[label=\Alph*)]
\item $\sqrt{5}mg$
\item $4mg$
\item $\sqrt{7}mg$
\item $3mg$
\end{enumerate}

\item Pastga tushayotgan lift $200 cm/s^2$ tezlanish bilan tormozlanmoqda. Lift ichida turgan 60 kg massali odamning lift poliga beradigan kuchi qanday bo‘ladi?
\begin{enumerate}[label=\Alph*)]
\item 480
\item 600
\item 120
\item 720
\end{enumerate}

\item Kosmik kema $20 m/s^2$ tezlanish bilan tik yuqoriga ko‘tarilmoqda. Bunda massasi 80 kg bo‘lgan kosmonavtning yuklanishi nechaga teng?
\begin{enumerate}[label=\Alph*)]
\item 2
\item 3
\item 1
\item 4
\end{enumerate}

\item O‘lik sirtmoq bo‘ylab uchayotgan samolyotning sirtmoqning eng pastki nuqtasidagi og‘irligi (vazni) qanday ifodalanadi?
\begin{enumerate}[label=\Alph*)]
\item $mg + \frac{m\vartheta^2}{R}$
\item $mg - \frac{m\vartheta^2}{R}$
\item $mg + \frac{m\vartheta^2}{2}$
\item $\frac{m\vartheta^2}{R}$
\end{enumerate}

\item Massasi 2000 kg bo‘lgan avtomobil radiusi 40 m bo‘lgan qavariq ko‘prikdan 36 km/soat tezlikda o‘tmoqda. Ko‘prikning o‘rtasida avtomobil ko‘prikka qancha kuch bilan bosadi (kN)?
\begin{enumerate}[label=\Alph*)]
\item 10
\item 30
\item 45
\item 15
\end{enumerate}

\item Massasi 1,0 t bo‘lgan avtomobil radiusi 100 m bo‘lgan egri yo‘lda harakatlanmoqda. Avtomobilning tezligi 18 km/soat bo‘lgan holda markazga intilma kuchni toping (N).
\begin{enumerate}[label=\Alph*)]
\item 180
\item 200
\item 250
\item 150
\end{enumerate}

\item 1000 kg massali avtomobil 50 m egrilik radiusiga ega bo‘lgan qavariq ko‘prikdan 36 km/soat tezlik bilan o‘tmoqda. Avtomobil ko‘prik o‘rtasiga kelganda unga qanday kuch bilan bosim ko‘rsatadi (N)? Ko‘prikning yuqori nuqtasida avtomobil unga bosim ko‘rsatmasligi uchun u qanday eng kichik tezlik bilan harakat qilishi kerak (km/soat)?
\begin{enumerate}[label=\Alph*)]
\item 1700 ; 100
\item 7800 ; 40
\item 7800 ; 80
\item 2400 ; 80
\end{enumerate}

\item 2000 kg massali avtomashina 36 km/soat tezlik bilan botiq ko‘prikdan o‘tmoqda. Ko‘prikning egrilik radiusi 100 m. Avtomashina ko‘prikning o‘rtasiga kelganda unga qanday kuch bilan bosim ko‘rsatadi (N)?
\begin{enumerate}[label=\Alph*)]
\item 31500
\item 21600
\item 5400
\item 6500
\end{enumerate}

\end{enumerate}


\begin{table}[h!]
\centering
\caption{18-\S. Og'irlik va tezlanishli harakat bo'yicha kalitlar}
\begin{tabular}{|l|c|c|c|c|c|c|c|c|c|c|}
\hline
\textbf{Test №} & 1 & 2 & 3 & 4 & 5 & 6 (13) & 7 (14) & 8 (15) & 9 (16) & 10 (17) \\ \hline
\textbf{Javob}  & A & A & D & A & C & D & D & C & B & B \\ \hline
\hline
\textbf{Test №} & 11 (19) & 12 (20) & 13 (21) & 14 (22) & 15 (27) & 16 (28) & 17 (29) & 18 (30) & 19 (31) & 20 (41) \\ \hline
\textbf{Javob}  & A & B & A & B & A & C & D & B & A & D \\ \hline
\hline
\textbf{Test №} & 21 (42) & 22 (43) & 23 (44) & & & & & & & \\ \hline
\textbf{Javob}  & C & C & B & & & & & & & \\ \hline
\end{tabular}
\end{table}